\documentclass[11pt]{article}
\usepackage{icmlko}
\begin{document}

\title{내 노트의 ``남은 것'' 표를 감사했다 --- 도메인 추세는 하나도 없었다}
\author{월드모델 실험실}
\date{노트 371 · 2026년 8월 1일}
\maketitle

\begin{abstract}
노트 370이 ``남은 것''에 한 줄을 적었다: \emph{아이돌 0.1375 --- 이번
창에서 계속 내려왔다($0.4860 \to 0.1375$). 왜.} 재 봤다. \textbf{내려오지
않았다.} 그 사이 두 처치를 \emph{짝}으로 잰 값이 $+0.0026$(6/12) ·
$+0.0114$(8/12)로 둘 다 양수이고, 다섯 시점의 95\% 구간이 서로 전부
겹친다($[0.052,0.817]$ $\cdots$ $[-0.147,0.422]$). \textbf{``계속
내려왔다''는 한 번 적합의 점추정을 노트 사이로 이어 읽은 것이고, 그것은
노트 358이 ``열 중 넷에서 부호가 틀린다''로 금지한 바로 그 읽기다.}

그래서 전부 감사했다. 이 창의 챔피언 넷에서 열 도메인의 점추정이 움직인
폭을 제 구간폭으로 나누면 \textbf{제일 큰 것이 0.41}(시장팝업)이고 중앙이
\textbf{0.07} 이다. \textbf{하나도 제 구간을 못 넘는다 --- 이 창에 도메인
추세는 없었다.} 진짜로 잰 변화는 전부 \emph{한 실험 안에서 짝으로} 잰
것들이었고(노트 354의 시장팝업 $-0.1577$ 등), 그것들은 유효하다.
\textbf{유효한 것과 아닌 것의 경계는 ``짝이냐''가 아니라 ``같은 실험
안이냐''다.}
\end{abstract}

\section{한 줄을 확인했다}

\begin{center}
\begin{tabular}{lrl}
\hline
시점 & 아이돌 & 95\% 구간 \\
\hline
353 (유보 25행) & 0.4860 & $[0.052,\,0.817]$ \\
353 (유보 51행) & 0.2805 & $[0.012,\,0.553]$ \\
354 & 0.2137 & $[-0.078,\,0.505]$ \\
358 & 0.1137 & $[-0.166,\,0.394]$ \\
370 & 0.1375 & $[-0.147,\,0.422]$ \\
\hline
\end{tabular}
\end{center}

\textbf{다섯이 서로 다 겹친다.} 그리고 그 사이 두 처치를 짝으로 잰 값은

\begin{center}
\begin{tabular}{lrrr}
\hline
 & 차 & SD & 부호 \\
\hline
358 팝업 등급 풀기 & $+0.0026$ & 0.0606 & 6/12 \\
370 채움 재현 & $+0.0114$ & 0.0277 & 8/12 \\
\hline
\end{tabular}
\end{center}

\textbf{둘 다 양수다.} 유일하게 진짜인 하락은 노트 353의 유보 넓힘
($0.4860 \to 0.2805$)인데, 그것은 하락이 아니라 \textbf{정정}이다 ---
25행 추정이 낙관이었다고 그 글이 이미 적었다.

\section{전부 감사했다}

\begin{center}\includegraphics[width=\linewidth]{figs/audit.pdf}\end{center}

\begin{center}
\begin{tabular}{lrrrrrr}
\hline
도메인 & 353 & 354 & 358 & 370 & 움직인 폭 & $\div$ 구간폭 \\
\hline
시장팝업 & 0.2761 & 0.1365 & 0.1715 & 0.1616 & 0.1396 & \textbf{0.41} \\
아이돌 & 0.2805 & 0.2137 & 0.1137 & 0.1375 & 0.1668 & 0.29 \\
웹툰 & 0.4508 & 0.4586 & 0.4712 & 0.4724 & 0.0216 & 0.17 \\
도서 & 0.4002 & 0.4063 & 0.4349 & 0.4236 & 0.0347 & 0.13 \\
팝업 & 0.4781 & 0.4905 & 0.4810 & 0.4575 & 0.0330 & 0.08 \\
펀딩 & 0.4106 & 0.3969 & 0.4120 & 0.4195 & 0.0226 & 0.06 \\
모바일 & 0.5967 & 0.5990 & 0.6014 & 0.6020 & 0.0053 & 0.04 \\
게임 & 0.5961 & 0.5965 & 0.5903 & 0.5911 & 0.0062 & 0.03 \\
애니 · 세계애니 & \multicolumn{4}{c}{$\cdots$} & 0.0018 & 0.01 \\
\hline
\end{tabular}
\end{center}

\textbf{하나도 1.0 을 못 넘는다.} 중앙이 0.07 이다. \textbf{이 창에서
도메인 점수가 ``추세''를 보인 적은 없다} --- 챔피언이 네 번 바뀌는 동안
움직인 것은 전부 제 잡음 안이다.

\section{그러면 무엇이 진짜였나}

\textbf{같은 실험 안에서 짝으로 잰 것들이다.} 이 창에서 제일 큰 도메인
변화들 --- 노트 354의 시장팝업 $-0.1577$(entry\_friction 막기), 노트 364의
팝업 $-0.0993$(0/12, 세계애니 학습 빼기), 노트 361의 아이돌 $+0.0587$
(10/12, 만화 빼기) --- 은 전부 \emph{한 실행 안에서 두 팔을 나란히} 잰
것이고, 유효하다.

\textbf{경계는 ``짝이냐''가 아니라 ``같은 실험 안이냐''다.} 노트 358이
``도메인 칸은 짝으로만 적는다''로 적었는데 그것만으로는 부족했다 ---
짝으로 잰 값이라도 \emph{다른 실험} 것끼리 빼면 안 된다. 노트 274 · 275가
``저장된 도메인 점수는 실행을 가로질러 빼면 안 된다''로 이미 적어 둔
것이고(지문/`store.comparable`), \textbf{내가 내 ``남은 것'' 표에서 그
규칙을 어겼다.}

\section{규약에 더하는 줄}

\begin{itemize}
\item \textbf{``남은 것'' 표에 도메인 추세를 적지 않는다.} 적으려면 그
추세 자체를 한 실험 안에서 짝으로 재고 적는다.
\item \textbf{수를 노트 사이로 이어 읽지 않는다} --- 노트 274의 지문
규칙을 산문에도 건다.
\item \textbf{자기 노트도 감사 대상이다.} 이번 것은 노트 370의 마지막
줄이 하루도 못 가서 틀린 경우다.
\end{itemize}

\section{남은 것}

\begin{center}
\begin{tabular}{ll}
\hline
갈래 & 지금 아는 것 \\
\hline
만화 KR/CN 583행 & 노트 79가 근거를 대고 뺐다. 지금 판에서 다시 재나 \\
세계애니 눈금 & 합친 백분위가 세계애니를 12/12 로 올렸다(노트 369) \\
팝업 실측 41행 & 노트 360 · 364 · 366 \\
\hline
\end{tabular}
\end{center}

\textbf{아이돌 줄은 지웠다} --- 물음이 아니었다.

\end{document}
